\begin{problem}
    Para calcular el monto final de dinero que se tendrá por medio de un interés simple mensual del $i\%$, después de una cantidad $n$ de meses, se utiliza la siguiente expresión:

    \[
    C_f = C_i \left(1 + \frac{i \cdot n}{100}\right),
    \]

    tal que $C_f$ es el monto final de dinero y $C_i$ es el monto inicial de dinero.

    Considera una inversión con un monto inicial de \SI{100000}[\$]{} al que se le aplica un interés simple mensual del \SI{20}{\%}.

    ¿Cuánto dinero se tendrá transcurridos \num{5} meses?

    \begin{enumerate}[label=\Alph*.]
        \item \SI{102000}[\$]{}
        \item \SI{110000}[\$]{}
        \item \SI{120000}[\$]{}
        \item \SI{200000}[\$]{}
    \end{enumerate}
\end{problem}
